\documentclass[12pt,a4paper]{article}
\usepackage[utf8]{inputenc}
\usepackage[T1]{fontenc}
\usepackage[ngerman]{babel}
\usepackage[margin=2cm,top=3.5cm,bottom=2.5cm]{geometry}
\usepackage{enumitem}
\usepackage{setspace}
\usepackage{parskip}
\usepackage{fancyhdr}
\usepackage{ragged2e}
\usepackage{tikz}
\usepackage{eurosym}


\title{}
\author{}
\date{}

\setlength{\parindent}{0pt}
\setlength{\parskip}{6pt}
\renewcommand{\baselinestretch}{1.15}

\newcommand\Committee{SEDE}
\newcommand\Fraktion{LEER}

\newcommand\Thema{Armee}
\newcommand\city{Ulm}
\newcommand\datum{11.05.2026}
\newcommand\timeVorb{07:30-09:00}
\newcommand\timeEinf{09:00-09:30}
\newcommand\timeFrakOne{09:30-11:00}
\newcommand\timePauseOne{11:00-11:15}
\newcommand\timeAuss{11:15-12:30}
\newcommand\timeMittag{12:30-13:00}
\newcommand\timeFrakTwo{13:00-13:30}
\newcommand\timeFinVerh{13:30-14:00}
\newcommand\timePlenar{14:00-15:15}
\newcommand\timeDebr{15:15-15:30}
\newcommand\politiker{Andrea Wechsler}
\newcommand\politikerOffice{Mitglied des Europäischen Parlaments}
\newcommand\stadtvertreter{Nalan Schmidt}
\newcommand\stadtvertreterOffice{Anlauf- und Koordinierungsstelle Bürgerdialog (Open Government) der Stadt Ulm}
\newcommand\localSupport{das Europe Direct Ulm}
\newcommand\sponsor{die Vertretung der Europäischen Kommission in München}
\newcommand\jefvorsitz{Farras Fathi}
\newcommand\gendervorsitz{Landesvorsitzender}
\newcommand\evpLeader{Tobias}
\newcommand\evpRoom{Kleiner Sitzungssaal}
\newcommand\sdLeader{Antonia}
\newcommand\sdRoom{Veranstaltungsraum Europe Direct}
\newcommand\reLeader{Alica}
\newcommand\reRoom{Besprechungszimmer Donaubüro I}
\newcommand\fifthLeader{Farras}
\newcommand\fifthRoom{Besprechungszimmer Öffentlichkeitsarbeit}
\newcommand\pfeLeader{Hannah}
\newcommand\pfeRoom{Besprechungszimmer BM III}
\newcommand\numSuS{27}
\newcommand\location{im Ulmer Rathaus}
\newcommand\fifthGroup{Linke}
\newcommand\anzahlcomm{3}
\newcommand\totcoms{26}
\newcommand\minmems{25}
\newcommand\maxmems{90}

\newcommand{\LogoBasePath}{../Logos}
% Shared logo helpers for SimEP LaTeX documents.

\providecommand{\LogoBasePath}{Logos}
\providecommand{\FraktionsLogoBasePath}{Folien/Bilder}

\providecommand{\constrainedlogo}[4][2cm]{%
  \raisebox{#3}{\includegraphics[width=#2,height=#1,keepaspectratio]{#4}}%
}

\providecommand{\jeflogo}{\constrainedlogo[2.5cm]{5cm}{0pt}{\LogoBasePath/jef.png}}
\providecommand{\simeplogo}{\constrainedlogo[2.5cm]{5cm}{0pt}{\LogoBasePath/simep.png}}
\providecommand{\financelogo}{\constrainedlogo[1.5cm]{4cm}{0pt}{\LogoBasePath/finance\_\city.png}}
\providecommand{\citylogo}{\constrainedlogo[1.5cm]{4cm}{0pt}{\LogoBasePath/\city.png}}
\providecommand{\fraktionslogo}{\constrainedlogo[2.5cm]{5cm}{0pt}{\FraktionsLogoBasePath/\Fraktion.png}}

\providecommand{\PlaceLogosOne}{%
\begin{tikzpicture}[remember picture, overlay]
        % Top left corner
        \node at ([xshift=10mm, yshift=-8mm] current page.north west) [anchor=north west, inner sep=0pt] {
            {\jeflogo}
        };

        % Top right corner
        \node at ([xshift=-7.5mm, yshift=-6mm] current page.north east) [anchor=north east, inner sep=0pt] {
            {\simeplogo}
        };


        % Bottom left corner
        \node[anchor=south west, inner sep=0pt, text width=4.5cm, align=left] at ([xshift=10mm, yshift=8mm] current page.south west) {
            \small Auftraggeber \& \newline Kooperationspartner:\\[3mm]
            \financelogo
        };


        % Bottom right corner
        \node[anchor=south east, inner sep=0pt, text width=4.5cm, align=right] at ([xshift=-10mm, yshift=10mm] current page.south east) {
            \small Unterst\"utzt durch:\\[3mm]
            \citylogo
        };

\end{tikzpicture}
}

\providecommand{\PlaceLogosRest}{%
\begin{tikzpicture}[remember picture, overlay]
        % Top left corner
        \node at ([xshift=10mm, yshift=-8mm] current page.north west) [anchor=north west, inner sep=0pt] {
            {\jeflogo}
        };

        % Top right corner
        \node at ([xshift=-7.5mm, yshift=-6mm] current page.north east) [anchor=north east, inner sep=0pt] {
            {\simeplogo}
        };

\end{tikzpicture}
}

\providecommand{\PlaceLogosFraktion}{%
\begin{tikzpicture}[remember picture, overlay]
        % Top left corner
        \node at ([xshift=10mm, yshift=-8mm] current page.north west) [anchor=north west, inner sep=0pt] {
            {\jeflogo}
        };

        % Top right corner
        \node at ([xshift=-7.5mm, yshift=-6mm] current page.north east) [anchor=north east, inner sep=0pt] {
            {\simeplogo}
        };

        % Upper-right caucus logo
        \node at ([xshift=-12mm, yshift=-35mm] current page.north east) [anchor=north east, inner sep=0pt] {
            {\fraktionslogo}
        };

        % Bottom left corner
        \node[anchor=south west, inner sep=0pt, text width=4.5cm, align=left] at ([xshift=10mm, yshift=8mm] current page.south west) {
            \small Auftraggeber \& \newline Kooperationspartner:\\[3mm]
            \financelogo
        };


        % Bottom right corner
        \node[anchor=south east, inner sep=0pt, text width=4.5cm, align=right] at ([xshift=-10mm, yshift=10mm] current page.south east) {
            \small Unterst\"utzt durch:\\[3mm]
            \citylogo
        };

\end{tikzpicture}
}

\begin{document}

\thispagestyle{empty}

\vspace*{2.5cm}

\begin{center}
\Large
Vorschlag für eine

\vspace{0.8cm}

\textbf{\large Verordnung des Europäischen Parlaments und des Rates\\[0.3em]
für ein Europäisches Klimagesetz}

\vspace{2cm}

\normalsize
Federführende Ausschüsse:

\vspace{0.5cm}

\textbf{Haushaltsausschuss (BUDG)}

\textbf{Ausschuss für Industrie, Forschung und Energie (ITRE)}

\textbf{Ausschuss für Verkehr und Tourismus (TRAN)}

\textbf{Ausschuss für Landwirtschaft und ländliche Entwicklung (AGRI)}
\end{center}


\PlaceLogosOne

\newpage

\begin{center}
\textbf{ENTWURF EINER ENTSCHLIEßUNG DES EUROPÄISCHEN PARLAMENTS}
\end{center}
zu dem Vorschlag der Kommission für eine Verordnung des Europäischen Parlaments und des Rates zur Realisierung von Klimaneutralität und der Einführung eines Europäischen Klimagesetzes in der Europäischen Union.

\vspace{-0.2cm}

\subsection*{DAS EUROPÄISCHE PARLAMENT,}

\begin{itemize}[leftmargin=2em, itemsep=0.5pt]
  \item gestützt auf Artikel 192(1) des Vertrags über die Arbeitsweise der Europäischen Union, auf deren Grundlage ihm der Vorschlag der Kommission unterbreitet wurde,
  \item gestützt auf seine Geschäftsordnung,
  \item gestützt auf die Mitteilung COM (2019) 640 der Kommission an das Europäische Parlament, den Europäischen Rat, den Rat, den Europäischen Wirtschafts- und Sozialausschuss und den Ausschuss der Regionen zum Europäischen Grünen Deal vom 11.12.2019,
  \item nach Stellungnahme des Europäischen Wirtschafts- und Sozialausschusses,
  \item nach Stellungnahme des Europäischen Ausschusses der Regionen,
  \item gestützt auf die Berichte der assoziierten Ausschüsse: Haushaltsausschuss (BUDG), Ausschuss für Industrie, Forschung und Energie (ITRE), Ausschuss für Verkehr und Tourismus (TRAN), Ausschuss für Landwirtschaft und ländliche Entwicklung (AGRI).
\end{itemize}

\subsubsection*{IN ERWÄGUNG NACHSTEHENDER GRÜNDE: (Alle Ausschüsse)}

\begin{enumerate}[label=(\arabic*), leftmargin=2em, itemsep=0.5pt]
  \item Wissenschaftliche Grundlage einer gemeinsamen Europäischen Klimapolitik bildet der Sonderbericht des Weltklimarats (IPCC), der die Notwendigkeit von Maßnahmen zur Begrenzung der globalen Erwärmung auf 1,5 °C darlegt.
  \item Als Teil des Kyoto-Protokolls ist es Ziel der Europäischen Union, den Vereinbarungen des Übereinkommens von Paris gerecht zu werden. Zentral ist dabei, dass der globale Temperaturanstieg deutlich unter 2\textdegree C bleiben und möglichst auf 1,5\textdegree C begrenzt werden soll. Ein gemeinsames Europäisches Klimagesetz ist hierfür ein essentieller Schritt.
  \item Für ergriffene Maßnahmen sollten der Grundsatz der Solidarität zwischen den Mitgliedstaaten gelten. Außerdem sollten das Wohlergehen der Bürger*innen, der Wohlstand der Gesellschaft, die Wettbewerbsfähigkeit der Wirtschaft sowie sichere Energie- und Lebensmittelversorgung zu erschwinglichen Preisen sichergestellt sein.
  \item Erst durch ein gemeinsames Vorgehen auf EU-Ebene können grenzüberschreitende Maßnahmen ergriffen werden, um dem Klimawandel als von Natur aus grenzüberschreitenden Prozess entgegenzuwirken. In Einklang mit dem Subsidiaritätsprinzip nach Artikel 5 über des Vertrages über die Europäische Union ist ein einheitlicher Rechtsrahmen daher effektiver als Maßnahmen auf Ebene der Mitgliedstaaten.
\end{enumerate}

\newpage

\noindent\textbf{HABEN FOLGENDE VERORDNUNG ERLASSEN:}

\subsection*{Artikel 1 – Gegenstand und Ziele}

\begin{enumerate}[leftmargin=2em, itemsep=2pt]
  \item Die Mitgliedstaaten der Union verpflichten sich zur Senkung von Treibhausgasemissionen auf null netto (Klimaneutralität) bis spätestens 2050 um zur Eindämmung der globalen Erwärmung auf unter 2\textdegree C, wenn möglich auf 1,5\textdegree C gegenüber vorindustrieller Zeit, beizutragen. (Alle Ausschüsse)
  \item Die zuständigen Organe auf Unions- und nationaler Ebene ergreifen, unter Berücksichtigung der Solidarität zwischen den Mitgliedstaaten, die notwendigen Maßnahmen, um die hier festgelegten Ziele zu erreichen. (Alle Ausschüsse)
\end{enumerate}

\subsection*{Artikel 2 – Geltungsbereich (Alle Ausschüsse)}

Die vorliegende Verordnung gilt für alle Maßnahmen der Klimapolitik, die im Hoheitsgebiet der Mitgliedstaaten der EU realisiert werden. Sie gibt das verbindliche Ziel vor, die im Übereinkommen von Paris festgelegten Temperaturziele sowie Klimaneutralität bis 2050 zu erreichen.

\subsection*{Artikel 3 – Finanzierung (BUDG)}

\begin{enumerate}[leftmargin=2em, itemsep=2pt]
  \item Für die Finanzierung des European Green Deal wird ein Investitionsfonds in Höhe von \euro{}1 Billion für das nächste Jahrzehnt bereitgestellt. Die Mitgliedstaaten beteiligen sich anteilig an der Finanzierung. Der jährliche Beitrag wird anhand folgender Formel berechnet: (Bruttoinlandsprodukt des Landes in \euro{} x 0,0075).
  \item Mitgliedstaaten, deren Wirtschaft einem besonderen und unverhältnismäßigen Druck ausgesetzt ist, können auf Antrag mit Mitteln des Investitionsfonds bei der Umstellung auf eine nachhaltige Wirtschaft unterstützt werden.
  \item Regionen, die in besonders starkem Ausmaß von der Umstellung auf eine nachhaltige Wirtschaft betroffen sind und viele Arbeitsplätze verlieren, sollen besonders gefördert werden. Für jeden verlorenen Arbeitsplatz werden der Region aus dem Investitionsfond \euro{}10.000 für Weiterbildungsmaßnahmen zur Verfügung gestellt.
\end{enumerate}

\subsection*{Artikel 4 – Industrie - Forschung - Energie (ITRE)}

\begin{enumerate}[leftmargin=2em, itemsep=2pt]
  \item Um CO2-arme Technologien nachhaltig zu unterstützen, wird für das EU-Emissionshandelssystem (EU-EHS) eine neue Klasse von kostenlosen Premium-Zertifikaten eingeführt. Diese Premium-Zertifikate werden ausschließlich an die Unternehmen, die CO2-arme Technologien verwenden, vergeben und besitzen den dreifachen Wert eines normalen Zertifikats im Rahmen des EU-EHS.
  \item Das Forschungsrahmenprogramm Horizon Europe soll die Entwicklung neuer, CO2-armer Technologien fördern. Dazu sollen jährlich mindestens 20\% des Gesamtetats in Höhe von \euro{}100 Milliarden in Forschungsvorhaben fließen, die einen direkten Bezug zum Klimaschutz haben.
  \item Der Mindestanteil an Erneuerbaren Energien (EE) im Energiemix der Mitgliedstaaten wird für die nächsten Jahrzehnte wie folgt festgelegt:
  \begin{enumerate}[label=\alph*., leftmargin=1.8em, itemsep=2pt]
    \item EE-Anteil bis 2030: 30\%
    \item EE-Anteil bis 2040: 50\%
    \item EE-Anteil bis 2050: 80\%
  \end{enumerate}
\end{enumerate}

\subsection*{Artikel 5 – Mobilität (TRAN)}

\begin{enumerate}[leftmargin=2em, itemsep=2pt]
  \item Der erlaubte durchschnittliche CO2-Ausstoß von neuzugelassenen Autos wird bis 2030 schrittweise von 95g/km auf 60g/km reduziert. Der vorgeschriebene Anteil von Elektroautos an den neuzugelassenen Autos soll bis 2030 auf 35\% steigen.
  \item Mit Mitteln des Investitionsfonds sollen zusätzliche Europäische Korridore für Hochgeschwindigkeitszüge geschaffen werden. Der Verkehr mit Hochgeschwindigkeitszügen in der EU soll im Vergleich zu 2015 bis 2030 verdoppelt und bis 2050 verdreifacht werden.
  \item Um Emissionen aus dem Flugverkehr stärker zu bepreisen, soll die Menge der kostenlosen Zertifikate für Airlines im EU-Emissionshandelssystem (EU-EHS) von gegenwärtig 82\% bis 2030 auf 75\% sinken. Künftig sollen auch für Flüge aus oder in die EU Zertifikate erforderlich sein.
\end{enumerate}

\subsection*{Artikel 6 – Landwirtschaft (AGRI)}

\begin{enumerate}[leftmargin=2em, itemsep=2pt]
  \item Die Gemeinsame Agrarpolitik (GAP) der EU soll mehr auf Nachhaltigkeit ausgerichtet werden. Dazu werden bis 2030 folgende Änderungen an der GAP umgesetzt:
  \begin{enumerate}[label=\alph*., leftmargin=1.8em, itemsep=4pt]
    \item Der Anteil der ökologisch bewirtschafteten Flächen an den landwirtschaftlichen Flächen soll von gegenwärtig 8\% auf 25\% steigen.
    \item Der Gebrauch von Pflanzenschutzmittel und Antibiotika soll um 50 \% gesenkt werden.
    \item Der Einsatz von Düngemittel soll um 20\% gesenkt werden.
  \end{enumerate}
  \item Im Rahmen der Absatzförderungspolitik sollen aus dem Investitionsfonds zusätzliche Mittel in Höhe von \euro{}40 Milliarden für die ökologische Landwirtschaft bereitgestellt werden.
  \item Um die Biodiversität in Europa zu schützen, werden bis 2030 10\% der gesamten Land- und Meeresgebieten unter strengen Naturschutz gestellt.
\end{enumerate}

\subsection*{Artikel 7 – Inkrafttreten (KEIN AUSSCHUSS)}

Diese Verordnung tritt am dritten Tag nach ihrer Veröffentlichung im Amtsblatt der Europäischen Union in Kraft. Diese Verordnung ist in allen ihren Teilen verbindlich und gilt unmittelbar in jedem Mitgliedstaat.

\end{document}